% ============================================================
%  08_conclusiones.tex
% ============================================================
\chapter{Conclusiones}

El desarrollo del proyecto \textbf{Gestión de Comandas} representa un esfuerzo
exitoso que logró cumplir con los objetivos planteados, consolidando una plataforma
integral para la modernización operativa de restaurantes pequeños y medianos. La
aplicación implementa funcionalidades clave como el registro de pedidos en tiempo
real, la vista de estado de mesas vía Supabase Realtime, la gestión diferenciada por
roles con Row Level Security, los reportes financieros calculados server-side mediante
funciones PostgreSQL, y la carta digital pública accesible sin app desde cualquier
dispositivo.

Este logro fue posible gracias a la adopción del modelo BaaS con Supabase, que
permitió eliminar la necesidad de un backend intermedio propio, reduciendo la
complejidad operacional y el costo de infraestructura. La elección de React Native
con Expo garantiza actualizaciones OTA sin pasar por las tiendas de aplicaciones,
lo que es especialmente valioso en un entorno de restaurante donde la continuidad
del servicio es crítica.

El trabajo en equipo bajo la metodología \textbf{Scrum} demostró ser fundamental
para el progreso organizado del proyecto, permitiendo una distribución efectiva de
responsabilidades y una supervisión constante del avance a través de sprints bien
definidos. La asignación clara de roles, combinada con el uso estratégico de
herramientas colaborativas como GitHub, Notion y Discord, facilitó la resolución
de desafíos técnicos y aseguró la entrega de incrementos funcionales en cada
iteración.

Si bien la plataforma alcanzó sus objetivos principales, algunas funcionalidades
avanzadas, como la integración con pasarelas de pago, el dashboard live global
para el administrador y las cancelaciones parciales, representan oportunidades de
mejora para futuras iteraciones que podrán construirse sobre la base tecnológica
y funcional sólida establecida en este desarrollo.

% ── 8.1 ──────────────────────────────────────────────────
\section{Recomendaciones}

\subsection{Futuros Desarrollos}

\begin{itemize}[leftmargin=1.5em]
  \item \textbf{Integración con pasarelas de pago:} Conectar con Mercado Pago,
        Stripe u otra pasarela disponible en el mercado local para automatizar
        la verificación de pagos QR y eliminar la confirmación manual.
  \item \textbf{Dashboard live global:} Implementar la vista global de operación
        del restaurante para el admin, mostrando el estado de todas las mesas
        y meseros en tiempo real.
  \item \textbf{Cancelaciones parciales y devoluciones:} Añadir un flujo
        dedicado para devolver o cancelar ítems individuales dentro de una orden.
  \item \textbf{Modo offline básico:} Implementar caché local de órdenes activas
        para operar sin conexión y sincronizar al reconectar, usando tecnologías
        como MMKV o SQLite en el cliente React Native.
  \item \textbf{Soporte multi-sucursal avanzado:} La tabla \texttt{restaurants}
        ya habilita multi-sucursal; la siguiente iteración puede añadir un panel
        de administración por cadena de restaurantes.
\end{itemize}

\subsection{Optimizaciones Técnicas}

\begin{itemize}[leftmargin=1.5em]
  \item \textbf{Migración al tier Pro de Supabase:} Eliminar el riesgo de pausa
        automática por inactividad antes del despliegue en producción real.
  \item \textbf{Índices adicionales:} Añadir índices compuestos en
        \texttt{orders(waiter\_id, status)} y
        \texttt{order\_items(order\_id, menu\_item\_id)} para optimizar
        las consultas más frecuentes en horas pico.
  \item \textbf{Pruebas automatizadas:} Incorporar pipelines de CI/CD con suites
        de testing unitario (Jest), integración y end-to-end (Detox) para reducir
        la deuda de regresión.
  \item \textbf{Monitoreo de rendimiento:} Integrar Sentry o similar para
        captura de errores y métricas de rendimiento en producción.
\end{itemize}

\subsection{Gestión y Comunicación}

\begin{itemize}[leftmargin=1.5em]
  \item \textbf{Documentación técnica continua:} Establecer un proceso sistemático
        para mantener actualizados el esquema de BD, las políticas RLS y las
        decisiones de arquitectura en el repositorio.
  \item \textbf{Protocolos de code review:} Definir checklists estandarizados
        para revisiones de código que aseguren calidad y consistencia.
  \item \textbf{Capacitación en herramientas:} Realizar sesiones de entrenamiento
        en el uso avanzado de Supabase (migraciones, funciones RPC, Realtime)
        para todo el equipo de desarrollo.
\end{itemize}